\documentclass{article}
\usepackage{graphicx} % Required for inserting images
\usepackage{amsmath}
\usepackage{listings}
\usepackage{xcolor}
\usepackage{amssymb}

\lstset{
    basicstyle=\ttfamily\small,
    breaklines=true,
    frame=single,
    language=Python
}

\title{Optimization}
\author{Chaithanya Thalapaneni, Abenezer Belay}
\date{October 2025}

\begin{document}

\maketitle

\section{Introduction}

\section{Simulated Annealing Algorithm}

The physical process of Annealing in metallurgy inspires Simulated Annealing. When you heat metal and then cool it slowly, the atoms have enough energy to move around and settle into a low-energy, stable configuration. If you cool it too quickly, the atoms get stuck in a suboptimal arrangement.

In optimization, we use this analogy. The "energy" of our system is the cost of a schedule—i.e., the number of constraint violations it has. We start with a random schedule and make small changes to it. Here is the key insight: early in the process (at high temperature), we accept many changes, even those that make the schedule slightly worse. Starting with high temperature prevents us from getting stuck in a local optimum—a schedule that's better than all its neighbors but not the best possible schedule overall. As we continue (lowering temperature), we become more selective, mainly accepting improvements.

The temperature parameter $\tau(t)$ controls this behavior. At high temperatures, we explore the solution space widely. At low temperatures, we exploit reasonable solutions by making minor refinements.

\subsection{Four Required Components}

To apply Simulated Annealing to the timetabling problem, we must define four components: the solution space, the cost function, the neighbor relation, and the optimization goal.

\subsubsection{Solution Space}

The solution space $C$ is the set of all possible complete assignments. A solution $S \in C$ is a complete assignment of all courses in $B$ to resources in $E$.

Let $B$ be the set of courses. Let $E = R \times T \times L$ where $R$ is the set of rooms, $T$ is the set of time slots, and $L$ is the set of lecturers.

A solution $S$ assigns to each course $b \in B$ a triple consisting of room, time slot, and lecturer. For example:
$$S(\text{Photogrammetric CV}) = (\text{SR\_A}, \text{Mon-09:15-10:45}, \text{Prof. Rodehorst})$$

This means Photogrammetric CV is scheduled in room SR\_A on Monday from 9:15 to 10:45, taught by Prof. Rodehorst.

\subsubsection{Cost Function}

The cost function $c : C \to \mathbb{R}$ measures how bad a schedule is. Lower cost means better schedule. We define:
$$c(S) = w_1 \cdot f_1(S) + w_2 \cdot f_2(S) + ... + w_n \cdot f_n(S)$$

where $f_i$ are individual penalty functions and $w_i$ are weights representing constraint importance.

For our implementation:
$$c(S) = 100 \cdot (\text{room conflicts}) + 100 \cdot (\text{lecturer conflicts}) + 50 \cdot (\text{workload}) + 10 \cdot (\text{gaps}) + 10 \cdot (\text{Friday})$$

For example, if a schedule has 2 room conflicts, 1 Friday afternoon class, and 3 student gaps:
$$c(S) = 100 \times 2 + 10 \times 1 + 10 \times 3 = 240$$

The penalties are classified by severity:
\begin{itemize}
    \item Hard constraints (weight = 100): Must be satisfied. Examples: no room conflicts, no lecturer conflicts
    \item Medium constraints (weight = 50): Important quality indicators. Example: maximum daily workload
    \item Soft constraints (weight = 10): Preferences. Examples: minimize gaps, avoid Friday afternoons
\end{itemize}

\subsubsection{Optimization Goal}

Our optimization problem is to find a solution $S^* \in C$ such that for all $S \in C$:
$$c(S^*) \leq c(S)$$

We call such $S^*$ an optimal solution.

\subsubsection{Neighbor Relation}

We define a binary relation $\sim$ on $C \times C$. For $S, S' \in C$, we say $S \sim S'$ if and only if we can obtain $S'$ from $S$ by exactly one of the following operations.

\textbf{Technique 1: Construct and destroy}

We select one course $b \in B$. Let $e$ be the current resource assigned to $b$ in $S$. We choose a new resource $e' \in E$ and create:
$$S' = (S \setminus \{(b,e)\}) \cup \{(b,e')\}$$

We remove exactly one assignment (the old one for course $b$) and add exactly one assignment (the new one for course $b$).

Example: Move Photogrammetric CV from (SR\_A, Mon-09:15, Prof. Rodehorst) to (SR\_H, Tue-09:15, Prof. Rodehorst).

\textbf{Technique 2: Swap Operation}

We select two courses $b_1, b_2 \in B$ with $b_1 \neq b_2$, and their resources $e_1, e_2 \in E$ where $(b_1,e_1) \in S$ and $(b_2,e_2) \in S$, then:
$$S' = (S \setminus \{(b_1,e_1), (b_2,e_2)\}) \cup \{(b_1,e_2), (b_2,e_1)\}$$

Example: Swap Virtual Reality at (SR\_H, Wed-13:30) with HCI Theory at (LH\_HK7, Mon-13:30).

We define the neighborhood of solution $S$ as:
$$N(S) := \{S' \in C | S \sim S'\}$$

\subsection{The Algorithm}

The Simulated Annealing algorithm proceeds as follows:

\begin{verbatim}
Algorithm: MinimizationSimulatedAnnealing(I)
Input: instance I
Output: solution S

1: let S be an initial solution S_0
2: let k be a constant
3: let t = 0
4: while we continue do
5:     let S' in N(S) be chosen randomly
6:     if c(S') <= c(S) then
7:         S := S'
8:     else
9:         with probability e^(-(c(S')-c(S))/(k*tau(t))) let S := S'
10:        otherwise let S unchanged
11:    end if
12:    t := t + 1
13: end while
14: Return(S)
\end{verbatim}

\subsection{Acceptance Probability}

When we generate a neighbor schedule $S'$ from current schedule $S$, we calculate the cost difference:
$$\Delta = c(S') - c(S)$$

If $\Delta < 0$, the neighbor is better—we always accept it. If $\Delta > 0$, the neighbor is worse—we accept it with probability:
$$P(\text{accept}) = e^{-\Delta / (k \cdot \tau(t))}$$

where $k$ is a constant and $\tau(t)$ is the current temperature.

Example: Current schedule has cost 200, neighbor has cost 250, temperature is 100, and $k=1$. Then $\Delta = 50$ and $P(\text{accept}) = e^{-0.5} \approx 0.606$ (60\% chance). This probability decreases as temperature drops.

Our implementation guarantees that hard constraints are never violated by validating each neighbor during generation, ensuring only valid neighbors are considered.
\end{document}